\section{变分态制备}
\label{sec:ansatz}

\subsection{优化目标}
\label{sec:points}

变分优化以两条扫描线 $\delta=-2$ 与 $\delta=0.85$ 上 $99$ 点 $s$ 网格（$s=0.01$--$0.99$）的每一点为目标：每点三相初态分别独立优化，记录三路最小值；目标哈密顿量覆盖图~\ref{fig:D04}(d) 中的平庸、拓扑与 AFM 区域。两条线上优化能量最低的初态分支由图~\ref{fig:D08D09}(a,b) 的背景色带标出。

\subsection{变分电路结构}

我们的变分电路遵循变分本征求解器范式~\cite{peruzzo2014variational}：
制备初态，再作用参数化拟设。经典循环最小化
\begin{equation}
E(\theta)=\langle\psi(\theta)|H|\psi(\theta)\rangle,
\end{equation}
其中 $|\psi(\theta)\rangle$ 为参数 $\theta$ 下的电路输出。
在每个目标参数点，三种参考初态分别在各深度 $p=1,2,3$ 独立优化。
先比较随机参数向量与零向量的能量，再对低能候选进行局域 L-BFGS 精修，
保留每个分支中能量最低的收敛结果，最后取三路最小能量 $E^*$ 及其参数。
保对称电路保持初态的对称性量子数。计算沿 $\delta=-2$ 与 $\delta=0.85$ 两条线逐点进行。
这类变分算法见综述~\cite{cerezo2021variational,tilly2022variational}，其中汇总代价函数设计、拟设族与 VQE 最佳实践。
理想态优化在经典计算机上进行，处理器执行所得电路以测量关联量。
对于处于同一扇区、且该扇区基态唯一的归一化试探态，谱分解给出
\begin{equation}
1-|\langle\psi_0|\psi(\theta)\rangle|^2
\leq \frac{E(\theta)-E_0}{\Delta_{\rm sec}}.
\label{eq:energy-fidelity}
\end{equation}
式~\eqref{eq:energy-fidelity} 通过扇区能隙 $\Delta_{\rm sec}$ 将试探态重叠与能量残差联系起来。

本文键旋转属于哈密顿量变分拟设：生成元取自相互作用项，转角独立优化~\cite{wecker2015progress}；该拟设族已有针对 XXZ 模型的研究~\cite{wiersema2020exploring}。
将对称性纳入态制备电路可把搜索限制在有关扇区~\cite{gard2020efficient}。
不同参考态提供不同的初始纠缠结构，这是非平庸态变分制备中的相关资源~\cite{ho2019efficient}。
三个参考态具有相同的有限链量子数，并提供不同的纠缠结构；各参数点的初态分支由其优化能量确定。

可训练性还依赖电路和代价函数结构。
趋近幺正设计的随机电路可具有指数缩小的梯度~\cite{mcclean2018barren}，局域与全局代价在浅层电路中也可呈现不同表现~\cite{cerezo2021cost}。
恒等块初始化~\cite{grant2019initialization} 和依据模型构造电路~\cite{wiersema2020exploring} 为结构化起点提供动机。
随深度出现的优化突变~\cite{campos2021abrupt}，以及特定噪声模型下的噪声诱导梯度压低~\cite{wang2021noise}，进一步刻画了电路深度和噪声在变分优化中的作用。

令 $M=L/2$。三个参考初态均在 $P=-2$ 扇区，其电路见图~\ref{fig:D06}；它们采用反射键轨道（orbit）拟设，镜面对映键对共享转角（奇键 $O_j\leftrightarrow O_{M+1-j}$，
偶键 $E_j\leftrightarrow E_{M-j}$）。拟设分两步构造：先用相匹配初态设定初始纠缠结构，再在所有键上施加由反射对称性约束的参数化旋转，参数量随键 orbit 数而非键数增长。
单层演化分奇偶两子层（$\alpha$ 为 orbit 编号，$\langle kl\rangle$ 为键）：
\begin{equation}
U_{o}=\prod_{\alpha}\prod_{\langle kl\rangle\in\mathcal{O}_{\alpha}}e^{-i\theta_{o,\alpha}(X_{k}X_{l}+Y_{k}Y_{l})/2}e^{-i\varphi_{o,\alpha}Z_{k}Z_{l}/2},
\end{equation}
\begin{equation}
U_{e}=\prod_{\alpha}\prod_{\langle kl\rangle\in\mathcal{E}_{\alpha}}e^{-i\theta_{e,\alpha}(X_{k}X_{l}+Y_{k}Y_{l})/2}e^{-i\varphi_{e,\alpha}Z_{k}Z_{l}/2}.
\end{equation}
同一键内 $XX{+}YY$ 与 $ZZ$ 对易，子层内各键不交因而对易；奇偶子层互不对易，必须保留两子层。每键实现为 $R_{XX}(\theta)R_{YY}(\theta)R_{ZZ}(\varphi)$ 门序列。每个 orbit 一组 $(\theta,\varphi)$，对本文 $L=8$ 的链，单层共 $L$ 个参数（$\delta=0$ 时 $\theta=\varphi$，减半）；拟设保持 $Z_{\rm tot}$、$\prod X$、$R$，始终位于 $P=-2$ 扇区。与初态配对重合的子层后作用：平庸与 AFM 初态偶键先，拓扑初态奇键先。
每相一个示意图见图~\ref{fig:D06}(a)--(c)：初态 + 所需拟设电路，分别为平庸、
拓扑、AFM 相。
理想电路下，试探态保持在目标对称性扇区。第~\ref{sec:comparison}~节比较实验与数值模拟的关联量。图~\ref{fig:D08D09}(a,b) 给出 $p=1$ 代价函数与精确基态能量的对比。

\subsection{与相关拟设的比较}
\label{sec:ansatz-comparison}

Yu 等从单态对出发，交替施加偶键与奇键 XXZ 旋转~\cite{yu2023simulating}；Wiersema 等的 XXZ 哈密顿量变分拟设采用类似的键分解~\cite{wiersema2020exploring}。
这些构造将所有同类奇键或偶键统一共享转角，使用 $4p$ 个参数，并在偶数格点开链上保持反射对称性。
本文仅约束镜像键共享转角，保留不同反射轨道的独立参数。
将同类奇偶轨道的转角设为相同值即可恢复统一共享拟设，因此放宽这些等角约束能在相同深度下扩大变分空间。

我们在两条扫描线的 $s=0.15,0.50,0.85$ 处，对 $p=1,2$ 和三种初态分别检验这一改善。
两种共享方式均比较零向量与 $64$ 个随机参数向量，在相同预算下用 L-BFGS 精修最低能的四个起点，再计算最低能收敛态相对精确基态的平方保真度 $F=|\langle\psi_0|\psi\rangle|^2$。
表~\ref{tab:ansatz-fidelity} 分别展示固定深度和固定参数数目下的改善。
轨道共享在两个测试深度下均提高了均匀链的保真度；在所列拓扑点，单层轨道拟设还优于相同参数数目的双层统一共享拟设。

\begin{table}[htbp]
\caption{两种共享方式的同预算保真度对照示例（$L=8$）。
$p$ 为变分层数，$N_{\rm par}$ 为独立转角数，$F$ 为最低能收敛态相对精确基态的平方保真度。
每组使用所列的相同初态。}
\label{tab:ansatz-fidelity}
\begin{ruledtabular}
\begin{tabular}{lrrr}
共享方式 & $p$ & $N_{\rm par}$ & $F$ \\
\hline
\multicolumn{4}{l}{$(s,\delta)=(0.50,-2.00)$，平庸初态} \\
统一共享 & $1$ & $4$ & $0.896314$ \\
反射轨道 & $1$ & $8$ & $0.931696$ \\
统一共享 & $2$ & $8$ & $0.986745$ \\
反射轨道 & $2$ & $16$ & $0.999997$ \\
\hline
\multicolumn{4}{l}{$(s,\delta)=(0.85,0.85)$，拓扑初态} \\
统一共享 & $2$ & $8$ & $0.925513$ \\
反射轨道 & $1$ & $8$ & $0.953804$
\end{tabular}
\end{ruledtabular}
\end{table}

Sun 等将目标纠缠结构的初始化与浅层 eSWAP 修饰分开，并讨论了键转角的反射约束~\cite{sun2023efficient}。
本文用三个参考态将这一初始化策略用于平庸、拓扑与 AFM 区域，并将 $XX{+}YY$ 与 $ZZ$ 转角分开以适应 XXZ 各向异性。
在 $\delta=0$ 时，两转角相等的键旋转与 eSWAP 仅差整体相位；对 $L=8$，反射共享将独立键层的 $7$ 个转角减为 $4$ 个。
理想修饰电路还保持完整的 $P=-2$ 扇区，而硬件高效拟设中的任意局域 Euler 旋转通常不保持这些自旋量子数~\cite{kandala2017hardware}。

两种共享方式的每个变分层均由两个不交键子层组成；采用显式 $R_{XX}R_{YY}R_{ZZ}$ 实现时，逻辑两比特旋转深度均为 $6p$。
完整硬件电路还包含初态制备、原生门分解和路由。

\begin{figure*}[!tp]
\centering
\includegraphics[width=\textwidth]{{../../figs/exp04/exp04_D06}}
\caption{三相初态与所需拟设电路：(a) 平庸初态；(b) 拓扑初态；
(c) AFM 初态。虚线框 $U^{(k)}$ 为第 $k$ 变分层；
$\theta^{(k)}_j$（$\varphi^{(k)}_j$）为第 $k$ 层第 $j$ 镜面对映键 orbit 共用的
$XX{+}YY$（$ZZ$）转角（奇键 $O_j\leftrightarrow O_{M+1-j}$，偶键
$E_j\leftrightarrow E_{M-j}$；$\delta=0$ 时 $\theta=\varphi$）。(a)(c)
偶键子层先作用，(b) 奇键子层先作用。
\label{fig:D06}}
\end{figure*}

\begin{figure*}[!tp]
\centering
\includegraphics[width=\textwidth]{{../../figs/exp06/exp06_D08D09}}
\caption{$p=1$ 代价函数，以及实验与数值模拟的对比。
(a)(b) $p=1$ 代价函数随 $s$ 变化：(a) $\delta=-2$，(b) $\delta=0.85$。
$E_0$（精确解，蓝色实线）与三路初态代价（平庸绿色虚线、拓扑橙色点线、
AFM 黑色点划线）；背景色带标出各网格点 $p=1$ 代价最低的初态分支，圆点标识各区间平均优化代价最低的分支。
(c)(d) 组合量 $Q$：(c) $\delta=-2$，(d) $\delta=0.85$。
数值模拟 $p=1$（黑）、$p=2$（灰）固定尺度曲线加实验散点
（$p=1$，每点 $16\,384$ shots；颜色按 (a)(b) 中最低能量初态分支分段）。误差棒为重复测量间的标准误差。
\label{fig:D08D09}}
\end{figure*}
